\chapter{Architectural Landscape of SVG Generation}
\label{ch:Landscape}

SVG generation is not a single problem with a single solution.
Some methods treat vector graphics as optimization targets and
adjust geometric parameters until a rendered image matches a
prompt or a reference image. Others learn from datasets and generate
SVGs through a trained decoder. A group of models avoids explicit
geometric primitives entirely and uses neural functions to
define shapes indirectly. 
These approaches reflect fundamentally different
assumptions about what a good SVG output should look like,
and they lead to very different trade-offs between visual
quality and structural organization.

This chapter introduces the thirteen models that form the whole thesis.
They were selected among a broader research landscape that also includes sequence-based
sketch models, retrieval-based vectorization systems, and large language model approaches to SVG
synthesis. That wider context is sketched below to
clarify where the selected models sit and why they were
chosen. The organizing principle is simple, these thirteen
models together cover the full range of design decisions of models
in two dimensions already mentioned: representation and supervision, that shape how SVG
structure emerges from a generation pipeline. A summary of
all thirteen models is provided in
Table~\ref{tab:model_overview}, and Figure~\ref{fig:taxonomy}
shows how they are grouped in this thesis.

\begin{figure}[tbp]
\centering
\begin{forest}
for tree={
  grow'=0, draw, rounded corners, fill=gray!10,
  font=\footnotesize, align=left,
  anchor=west, child anchor=west, parent anchor=east,
  l sep=5mm, s sep=1.2mm, inner sep=2pt,
  tier/.option=level,
  edge path={\noexpand\path[\forestoption{edge}]
    (!u.parent anchor) -- +(2.5mm,0) |- (.child anchor)\forestoption{edge label};}
}
[13 analyzed\\models
  [Optimization-based\\(per instance)
    [Reconstruction-supervised:\\LIVE{,} SAMVG]
    [Semantic-supervised{,} parametric:\\VectorFusion{,} SVGDreamer{,} SVGDreamer++]
    [Semantic-supervised{,} neural implicit:\\NiVeL{,} NeuralSVG]
  ]
  [Learning-based\\(trained{,} forward generation)
    [Sequential tokens: DeepSVG{,} DeepIcon]
    [Latent diffusion: SVGFusion]
    [Layered diffusion: LayerTracer]
  ]
  [Hybrid
    [Im2Vec: trained without vector supervision]
    [T2V-NPR: trained path prior\\+ per-prompt optimization]
  ]
]
\end{forest}
\caption{Taxonomy of the thirteen analyzed models by generation 
strategy, with representation and supervision as the second level.}
\label{fig:taxonomy}
\end{figure}

An early line of work treats vector drawings as ordered 
sequences of commands or strokes. Starting from 
SketchRNN (Section~\ref{sec:sequential_repr}), later 
methods extended this to full SVG command prediction, 
including 
Sketchformer~\cite{ribeiroSketchformerTransformerbasedRepresentation2020} 
and related sequence 
models~\cite{ganinSynthesizingProgramsImages2018, 
lopesLearnedRepresentationScalable2019b}. These approaches 
established that SVG structure could be learned 
statistically, but long command chains make training 
difficult and generalization to complex scenes remains 
limited by the scarcity of large-scale SVG datasets.

A second direction introduced differentiable rasterization as
a bridge between vector parameters and image-space supervision.
This opened the door to reconstruction-based
vectorization and, later, to semantically supervised generation
using pretrained diffusion models. Most of the optimization-based
models in this thesis belong to this lineage.

A third direction replaces explicit geometric primitives with
neural implicit functions. Instead of optimizing Bézier
coordinates directly, these methods run optimization in neural
parameter space, where an MLP defines shape boundaries
indirectly. This approach has connections to neural radiance
fields~\cite{mildenhallNeRFRepresentingScenes2020} and
represents an attempt to improve structural coherence through
representational design rather than supervision alone.

More recently, large language models and multimodal variants
have entered SVG generation. Methods such as
StarVector~\cite{rodriguezStarVectorGeneratingScalable2025},
OmniSVG~\cite{yangOmniSVGUnifiedScalable2025},
LLM4SVG~\cite{xingEmpoweringLLMsUnderstand2025}, and
InternSVG~\cite{wangInternSVGUnifiedSVG2025} treat SVG as
structured code and generate it through transformer
architectures pretrained on collection of large-scale texts and images. 
In one sense, SVG is a natural fit for language
models: the file format is sequential XML commands, not
unlike a programming language. Visual alignment improves
substantially compared to earlier sequence models, and
these approaches handle a wide range of generation and
editing tasks. However, they are not included in the main
analysis here. The core question of this thesis is how
specific architectural decisions impact the structure
of the output and that relationship is difficult to
trace through billions of parameters. A more detailed
justification is provided in Appendix~\ref{app:llm_exclusion}.

\section{Optimization-Based Methods}
\label{sec:optimization}

Optimization-based methods share a common logic: SVG
primitives are randomly or heuristically
initialized and then iteratively refined until the rendered
output satisfies a supervision signal which is either prompt or target image.
The SVG is the variable being optimized, not the output of a forward pass.
What varies across methods in this family is what that
supervision signal is and how geometry is represented
during optimization.

Almost all of them rely on DiffVG~\cite{liDifferentiableVectorGraphics2020a} to work. 
NiVeL is the exception: it optimizes implicit fields that are evaluated on a pixel grid 
and only turned into curves afterwards with marching squares, so DiffVG is not part of 
its optimization loop~\cite{thamizharasanNIVeLNeuralImplicit2024}.
The optimization loop across all methods 
follows the same render-compare-update cycle, 
but differs in what loss drives it and how geometry is represented.
Table~\ref{tab:arch_optbased} summarizes how each 
model in this group implements those differently.

\subsection{Reconstruction-Supervised}

LIVE~\cite{maLayerwiseImageVectorization2022} is the simplest form of optimization-based vectorization. 
It starts from an empty canvas, finds where the biggest visual gap is between the current SVG and the target image, 
drops a new path there, and optimizes everything together. It repeats this process layer by layer until the reconstruction is close enough. 
Because paths are added one at a time from coarse to fine, the output tends to have a layered structure where each path has a recognizable role. 
SAMVG~\cite{zhuSAMVGMultistageImage2023} follows the same optimization loop but starts from a better place: instead of finding gaps by 
pixel difference, it uses the Segment Anything Model~\cite{kirillovSegmentAnything2023} to identify semantically meaningful 
regions before any optimization begins. The loss function is the same. The difference in output quality comes almost entirely 
from where optimization starts, which suggests that initialization matters as much as what you optimize for.

Im2Vec~\cite{reddyIm2VecSynthesizingVector2021} does not fit cleanly into either group. 
It trains an encoder-decoder to predict vector primitives from a raster image in a single forward pass, 
using differentiable rendering and reconstruction loss during training. The training data are raster images only, 
the model never sees an SVG~\cite{reddyIm2VecSynthesizingVector2021}. At inference there is no iterative refinement. 
It is included here because its core contribution is the training objective, not a learned generative model over SVG structure.

\subsection{Semantic-Supervised}

When no target image exists, the optimization loop needs a different 
kind of signal. These methods use a pretrained diffusion model to judge whether the 
current SVG looks like what the text prompt describes. The mechanism 
behind this is Score Distillation Sampling 
(SDS)~\cite{pooleDreamFusionTextto3DUsing2022}, originally developed 
for text-to-3D generation and described in 
Section~\ref{sec:sds_vsd}: render the SVG, add noise, let the 
diffusion model predict that noise and 
use the difference between predicted and added noise as a gradient to update the SVG parameters.

VectorFusion~\cite{jainVectorFusionTexttoSVGAbstracting2022} was 
among the first to apply SDS to SVG generation, using Stable 
Diffusion~\cite{rombachHighResolutionImageSynthesis2021} as the 
supervision backbone (Figure~\ref{fig:sds_loop}). The outputs are 
semantically aligned with the prompt but tend to be over-smoothed, 
a known side effect of standard SDS. 
SVGDreamer~\cite{xingSVGDreamerTextGuided2024} addresses this by 
replacing SDS with Vectorized Particle-based Score Distillation 
(VPSD, Section~\ref{sec:sds_vsd}), which models the SVG as a 
distribution of particles rather than a single point estimate 
(Figure~\ref{fig:svgdreamer_vpsd}). SVGDreamer also introduces 
explicit control over how many primitives are used and how they 
are organized into semantic layers. 
SVGDreamer++~\cite{xingSVGDreamerAdvancingEditability2025} extends 
this further with a hierarchical refinement mechanism that can add 
or remove primitives during generation based on how much detail 
each region needs.


\begin{figure}[htbp]
\centering
\includegraphics[width=0.85\textwidth]{SDS.png}
\caption{SDS optimization loop in VectorFusion~\cite{jainVectorFusionTexttoSVGAbstracting2022}.}
\label{fig:sds_loop}
\end{figure}


\begin{figure}[htbp]
\centering
\includegraphics[width=0.85\textwidth]{VPSD!.png}
\caption{VPSD optimization loop in SVGDreamer~\cite{xingSVGDreamerTextGuided2024}.}
\label{fig:svgdreamer_vpsd}
\end{figure}

The two loops look alike, but they differ in what is trained. 
In the SDS loop (Figure~\ref{fig:sds_loop}), everything except the SVG is frozen. 
The rendered image is augmented, encoded into the latent space of Stable Diffusion 
and noised, and the frozen diffusion model predicts the noise. The weighted difference 
between predicted and added noise is the only signal, and it flows back through DiffVG 
into the paths of a single SVG. In the VPSD loop (Figure~\ref{fig:svgdreamer_vpsd}), 
the pretrained diffusion model is still frozen, but two parts are added. First, there 
are $k$ SVGs (particles) instead of one, and a LoRA copy of the diffusion model is 
fine-tuned on their noisy renders with its own loss. This copy learns what the current 
particles look like, and the update for the SVG parameters comes from the difference 
between the frozen model and the LoRA model, not from the added noise. Second, a reward 
model scores the rendered particles and reweights them, which the authors report gives 
better visual aesthetics and faster convergence~\cite{xingSVGDreamerTextGuided2024}. 
So in SDS only the SVG learns. In VPSD the particles and the LoRA learn together.

\subsection{Neural Implicit Representation}

NiVeL~\cite{thamizharasanNIVeLNeuralImplicit2024} and 
NeuralSVG~\cite{polaczekNeuralSVGImplicitRepresentation2025} take 
a different approach to the same optimization loop. 
They store geometry as the weights of a small MLP, instead of 
storing geometry as explicit B\'{e}zier coordinates that are 
adjusted directly. The network takes a number (shape index) 
and outputs the geometry for that shape. Nothing is stored as an explicit path.

This brings different properties to optimization based approaches. Compared to 
parametric methods, where each path is an independent object and 
optimization can push them in different directions. But with an 
MLP, all shapes come from the same learned function, which 
creates a bias toward global coherence. NiVeL uses this to 
decompose scenes into layered components, where each layer is 
a separate neural field whose zero-level set defines a shape 
boundary. NeuralSVG assigns each shape a fixed index and 
optimizes the network so that each index consistently maps to 
the same visual region across iterations. The outputs from 
both tend to be cleaner and more organized than purely 
parametric optimization, but inference is slower because the 
network has to be optimized from scratch for every new prompt.

\begin{table}[htbp]
\centering
\caption{Architectural overview of optimization-based 
SVG generation methods. Trainables: what is updated, 
per-instance parameters optimized at inference time, 
except for Im2Vec, whose network weights are trained once. 
Init.: how primitives are placed before optimization. 
Supervision: the loss that drives the update. 
Representation: how geometry is stored while it is optimized.}
\label{tab:arch_optbased}
\footnotesize
\begin{tabularx}{\textwidth}{@{}l l l l X@{}}
\toprule
\textbf{Model} & \textbf{Trainables} & \textbf{Init.} 
& \textbf{Supervision} & \textbf{Representation} \\
\midrule
LIVE         & SVG param.          & Error-driven     
             & Pixel recon.        & Parametric       \\
SAMVG        & SVG param.          & Mask-based       
             & Pixel recon.        & Parametric       \\
Im2Vec       & Encoder + RNN       & Image encoding   
             & Pixel recon.        & Parametric paths \\
VectorFusion & SVG param.          & Random/Trace     
             & SDS                 & Parametric       \\
SVGDreamer   & SVG param. + LoRA   & Attention (SIVE) 
             & VPSD                & Parametric       \\
SVGDreamer++ & SVG param. + LoRA   & Random + HIVE    
             & VPSD + HIVE         & Structured param.\\
NeuralSVG    & MLP weights + LoRA  & Saliency-based   
             & SDS                 & Neural implicit  \\
NiVeL        & MLP weights         & Random weights   
             & SDS                 & Neural implicit  \\
\bottomrule
\end{tabularx}
\end{table}

Table~\ref{tab:arch_optbased} summarizes the 
architectural components of the optimization-based 
methods. For all models in this group except Im2Vec, the 
``Trainables'' column refers to per-instance 
parameters optimized at inference time rather than 
learned model weights. Apart from Im2Vec, none of these 
methods requires a training dataset; supervision is 
provided either by the input image directly or by a 
frozen pretrained model. Im2Vec is the exception: its 
trainables are network weights, learned once from a 
dataset of raster images.



\section{Dataset-Driven Methods}
\label{sec:datasetdriven}

Dataset-driven methods learn structural patterns from collections of vector graphics
and generate through a trained forward pass. The
optimization happens during training, not at inference.
Outputs reflect statistical regularities in the training
data rather than the iterative refinement of any single
instance. These methods are more constrained by the training corpus
but can produce more consistent and visually appealing results
without the computational cost of optimization at inference.

DeepSVG~\cite{carlierDeepSVGHierarchicalGenerative2020}
is the foundational model in this family. It learns a
hierarchical latent space over SVG path commands using
a transformer-based encoder-decoder trained on the
SVG-Icons8 dataset.
The latent space separates path-level and command-level
representations, allowing the model to capture both coarse
shape organization and fine geometric detail. At inference,
SVG commands are generated autoregressively from a sampled
latent vector. DeepSVG demonstrated that structured SVG
generation from a learned latent space is feasible, and
also exposed the difficulty of maintaining geometric
validity across long command sequences.

DeepIcon~\cite{bingDeepIconHierarchicalNetwork2024} applies a similar framework
to icon vectorization specifically, trained on the same
SVG-Icons8 dataset of simple icons. The domain-specific training
produces outputs with more consistent style and geometry
than general-purpose models, though generalization beyond
the training distribution is limited by design.

LayerTracer~\cite{songLayerTracerCognitiveAlignedLayered2025} approaches the problem
differently. Rather than predicting SVG commands from a
latent vector, it models the construction process of
vector graphics layer by layer. A diffusion-based
architecture generates intermediate structural
representations that guide each new layer, approximating
the sequential way designers build graphics in practice.
This process-aware formulation produces outputs with more
meaningful layer organization than direct command
prediction, with direct consequences for editability.

T2V-NPR~\cite{zhangTexttoVectorGenerationNeural2024} combines dataset training
with per-prompt optimization. It first trains a neural path VAE
on the FIGR-8-SVG icon dataset, so the latent space learns
what a regular, smooth path looks like. For a new prompt there is
no forward pass. Path latents in this learned space are optimized
with VSD from a pretrained diffusion model, and the result is then
refined with layer-wise vectorization~\cite{zhangTexttoVectorGenerationNeural2024}.
So T2V-NPR is a hybrid: a trained path prior with per-instance
optimization on top. It is described in this section because its
path prior comes from a dataset, but its inference works like
the optimization methods above.

SVGFusion~\cite{xingSVGFusionScalableTexttoSVG2025} is the most architecturally
complex model in this family. It first learns a continuous
latent space of vector graphics using a Vector Path
Variational Autoencoder (VP-VAE), trained to compress path
sequences into dense latent representations. Generation
then occurs through a Vector Semantic Diffusion Transformer
(VS-DiT), which operates entirely in this learned latent
space. This makes SVGFusion distinct from the
semantic-supervised optimization methods in the previous
section: where VectorFusion and SVGDreamer use diffusion
gradients to guide per-instance optimization in pixel space,
SVGFusion uses a diffusion process over a learned vector
distribution to sample new SVGs directly. The outputs
reflect the statistical structure of the training corpus
rather than the convergence of any optimization loop. 
Table~\ref{tab:arch_dataset} summarizes the 
architectural components of the dataset-driven 
methods.

\begin{table}[htbp]
\centering
\caption{Architectural overview of dataset-driven 
SVG generation methods. Training data refers to 
the dataset used for offline model training.}
\label{tab:arch_dataset}
\footnotesize
\begin{tabularx}{\textwidth}{@{}l l l l X@{}}
\toprule
\textbf{Model} & \textbf{Trainables} & \textbf{Data} 
& \textbf{Supervision} & \textbf{Representation} \\
\midrule
DeepSVG      & VAE + Transformer   & SVG-Icons8       
             & Recon. + KL         & Sequential tokens \\
DeepIcon     & Transformer         & SVG-Icons8       
             & Reconstruction      & Sequential tokens \\
SVGFusion    & VP-VAE + VS-DiT     & SVGX             
             & Recon. + Denoising  & Continuous latent \\
T2V-NPR      & NPR-VAE + LoRA      & FIGR-8-SVG       
             & VSD                 & Neural path repr. \\
LayerTracer  & DiT + LoRA          & 20k layered SVGs 
             & Denoising           & Layered sequences \\
\bottomrule
\end{tabularx}
\end{table}

Unlike the optimization-based group, every model here has a 
training corpus. The representation column also shifts 
away from parametric paths toward sequential tokens, latent 
spaces, and layered sequences, reflecting a fundamentally 
different approach to how SVG structure is encoded.



Table~\ref{tab:model_overview} summarizes the thirteen models 
analyzed in this thesis. The \textit{Inference} column 
distinguishes between methods that optimize vector parameters 
at inference time (\textit{Iterative Opt.}) and methods that 
generate SVGs through a single trained forward pass 
(\textit{Forward Dec.}). The \textit{Supervision} column 
indicates the primary signal driving generation: 
\textit{Pixel recon.} refers to reconstruction against a 
target raster image without a pretrained model, \textit{Self-Sup.} to semantic 
signals derived from frozen pretrained models (diffusion or 
CLIP), and \textit{Supervised} to training against 
ground-truth SVG data. \textit{Hybrid} in the Inference 
column marks methods that combine a trained model with a 
per-instance optimization or vectorization step. Task types distinguish between 
\textit{Vectorization} (raster image as input, SVG as 
output), \textit{Text-to-SVG} (text prompt as input), and 
\textit{Unconditional Generation} (no input at inference --- SVG 
sampled directly from a learned latent space).

\begin{table}[htbp]
\centering
\caption{Overview of the thirteen models analyzed in this 
thesis, organized by generation strategy.}
\label{tab:model_overview}
\footnotesize
\begin{tabularx}{\textwidth}{l l l l l X}
\toprule
\textbf{Model} & \textbf{Task} & \textbf{Inference} & 
\textbf{Input} & \textbf{Output} & \textbf{Supervision} \\
\midrule
LIVE         & Vectorization      & Iterative Opt. & Image    & Paths        & Pixel recon. \\
SAMVG        & Vectorization      & Iterative Opt. & Image    & Paths        & Pixel recon. \\
Im2Vec       & Vectorization      & Forward Dec.   & Image    & Paths        & Pixel recon. \\
\midrule
VectorFusion & Text-to-SVG        & Iterative Opt. & Text     & Paths        & Self-Sup.  \\
SVGDreamer   & Text-to-SVG        & Iterative Opt. & Text     & Paths        & Self-Sup.  \\
SVGDreamer++ & Text-to-SVG        & Iterative Opt. & Text     & Paths        & Self-Sup.  \\
NiVeL        & Text-to-SVG        & Iterative Opt. & Text     & Layers/Paths & Self-Sup.  \\
NeuralSVG    & Text-to-SVG        & Iterative Opt. & Text     & Layers/Paths & Self-Sup.  \\
\midrule
DeepSVG      & Unconditional      & Forward Dec.   & None     & Paths        & Supervised \\
DeepIcon     & Vectorization      & Forward Dec.   & Image    & Layers/Paths & Supervised \\
SVGFusion    & Text-to-SVG        & Forward Dec.   & Text     & Layers/Paths & Supervised \\
T2V-NPR      & Text-to-SVG        & Hybrid         & Text     & Layers/Paths & Sup. + Self-Sup. \\
LayerTracer  & Text-to-SVG        & Hybrid         & Text/Img & Layers/Paths & Supervised \\
\bottomrule
\end{tabularx}
\end{table}

The split between Iterative Opt. and Forward Dec. in the 
Inference column maps almost cleanly onto the two families: 
every optimization-based model except Im2Vec refines at 
inference time, while every dataset-driven model except 
T2V-NPR generates through a forward pass. Im2Vec is trained 
and then decodes in a single pass. T2V-NPR trains a path 
prior and then optimizes for each prompt. LayerTracer is 
labeled Hybrid because it combines a trained diffusion model 
with a per-instance vectorization step.


The thirteen models described here differ in representation 
and supervision, and those two dimensions do not operate 
independently. A model's representation constrains what 
supervision signals are meaningful, and the choice of 
supervision often determines which representations are 
practical. The next two chapters trace these interactions: 
first through representation (Chapter~\ref{ch:representation}), 
then through supervision (Chapter~\ref{ch:supervision}).